\section{Introduction}
\label{sec:introduction}

OSEK applications are not just sequential C programs with a few library calls. They combine
task-level control flow, scheduler-visible operating-system APIs, event synchronization, and
interrupt service routines (ISRs) whose interference is both asynchronous and highly
domain-specific. The OSEK/VDX operating-system specification models a single-processor,
resource-aware, event-driven execution environment with explicit task states, scheduling
policies, and a standard API surface for services such as \texttt{ActivateTask},
\texttt{TerminateTask}, \texttt{SetEvent}, and \texttt{WaitEvent}~\cite{osek_spec}. Real
RTOS code bases such as Trampoline make heavy use of exactly this execution model and are
therefore poor fits for a verification flow that assumes ordinary function-call semantics or
unstructured asynchronous interference~\cite{trampoline_repo}.

CBMC already offers a strong baseline for C verification. Its normal architecture builds a
GOTO model from source code, symbolically executes that model, and discharges the resulting
equation with a SAT or SMT solver~\cite{cbmc_architecture,cbmc_repo}. However, a stock
CBMC workflow leaves two important gaps for OSEK-style software.

First, scheduler-facing OS calls are semantically richer than ordinary C functions. Treating
\texttt{ActivateTask} or \texttt{ChainTask} as plain calls either loses the scheduling meaning
entirely or forces users to encode scheduler behavior manually inside harness code. Second, a
naive interrupt model can place ISR interference at many irrelevant source locations. That
kind of over-approximation may still be sound, but it creates unnecessary state-space growth
and can produce failures that are artifacts of the modeling style rather than of the program.

The branch documented in this report addresses both problems with a deliberately split design.
An \texttt{os-api} layer plugs into \texttt{goto-symex} so that recognized OSEK calls affect
symbolic scheduler state directly. A separate \texttt{interleaving-analysis} library then
computes a manifest of plausible ISR insertion points from an existing GOTO model, and an
independent \texttt{aib} rewriter materializes only those guarded calls in a copied project
tree. The resulting workflow is more modular than embedding every policy into one stage of the
toolchain, and it exposes a concrete JSON contract between static interference analysis and
source-to-source instrumentation.

This report is an implementation-facing paper-style summary of that pipeline. It rewrites the
seed material in \texttt{docs/reports/Introduction.pdf} and
\texttt{docs/reports/Background.pdf}, but treats the checked-in code as the source of truth
when the prose and implementation differ. The report makes three concrete contributions:

\begin{itemize}[leftmargin=1.5em]
  \item It explains how OSEK-aware symbolic execution is integrated into CBMC through
  command-line options, dispatcher logic, and a generic scheduler action interface.
  \item It documents the manifest-driven ISR pipeline, including target-function discovery,
  read/write-set analysis, candidate-line reporting, and guarded source rewriting.
  \item It summarizes the available benchmark evidence, including cases where the improved
  pipeline is slower, cases where it is audited as the more faithful model, and one Trampoline
  walkthrough that shows how aggressively the targeted manifest shrinks the naive ISR
  injection space.
\end{itemize}

The rest of the report proceeds from context to mechanism to evidence. \Cref{sec:background}
briefly reviews the OSEK execution model, the stock CBMC flow, and the main families of
interrupt-aware model checking. \Cref{sec:method} describes the improved pipeline in detail.
\Cref{sec:evaluation} summarizes the evidence available in the checked-in benchmark artifacts,
and \Cref{sec:discussion} closes with practical limitations that matter for future
maintenance.
